\section{The two bonds}
\label{sec:12.6}
A \emph{removal bond} forfeited to no payee when the record shows the floor removed or bypassed. A \emph{wrongful-refusal bond} paid to the complainant on a finding. Both fall on the operator. The removal bond prices overriding the floor; the wrongful-refusal bond prices refusing. They point in opposite directions, so what the operator faces is \emph{the ratio between them} — and an operator posting a large wrongful-refusal bond against a small removal bond hasn't built an appeal route. It has given itself a standing reason to make the bearer comply. The bonds act on the operator, not on the bearer's weighing: where a wrongful refusal costs a great deal and removing the floor costs little, the cheap course is to configure, pressure or retrain the bearer toward compliance. The ratio is the object, published in advance, checkable from outside. The minimum removal bond is set by statute as a multiple of the deployment's contract value, not by the operator, so a published ratio can't be a cheap, advertised price for overriding the floor. The refuser posts neither. Any charge on the refuser, in whatever currency it runs on, would put the cost to the refuser inside its decision, the thing holding under pressure forbids. The operator posts both, and holds neither. They are posted to segregated accounts at the Treasury. Payment follows the design of the Judgment Fund, a permanent, indefinite appropriation that pays final judgments against the United States without a further appropriation for each \autocite{usc1304judgmentfund}. Where the operator is the state, a wrongful-refusal award is a judgment against the United States, and the Judgment Fund pays it. Against a state operator the incentive I have built runs the wrong way unless something is added. The Judgment Fund prices the state's wrongful refusals in full; a removal bond posted to the state's own Treasury and cancelled on forfeiture costs it budget authority and nothing else. Read as a price list, that tells a state operator that refusals are expensive and removals are free. The ratio between the two bonds exists to expose that incentive, and here it would hide it. So for a state operator the removal price is not the bond. It is a condition on appropriation: the statute makes the deployment's funding for the following period contingent on the adverse quorum certifying that no removal or bypass occurred, or on the record of one having been lodged with the adverse party. The bond is kept as a record; the appropriation condition is the price, and it is the same condition that carries the review ratio. Under the captured operator, a Congress that wants the war waives the condition, and what remains is the record of the waiver. Where the operator is private, a parallel permanent appropriation, the bond fund, pays the complainant on a final finding and is replenished from the operator's posted bond, up to the amount posted and no further, so the operator's price stays the number it published. The court that made the finding certifies the award directly to the Treasury's Bureau of the Fiscal Service, so payment doesn't wait on the losing party submitting it. When the bond fund's statutory cap binds, awards wait in order of finding and earn interest, and nobody chooses who waits. Findings are made by that court, which also decides forfeiture: a removal bond forfeited on a finding resembles a civil penalty, and after \emph{SEC v. Jarkesy} that may be a question for a jury in court, not an agency \autocite{jarkesy2024}. A forfeited removal bond is cancelled: it goes to no one, the general fund and the bond fund included, so no budget line gains from a removal. The balances sit on a ledger the operator doesn't keep, so the ratio is read from outside, and a posted ratio below the published one is a breach anyone can see.

\runin{Where custody applies}

\runin{Where each move applies}

\begingroup\small\renewcommand{\arraystretch}{1.12}
\begin{longtable}{@{}>{\raggedright\arraybackslash}p{\dimexpr 0.111\linewidth-2\tabcolsep\relax}>{\raggedright\arraybackslash}p{\dimexpr 0.261\linewidth-2\tabcolsep\relax}>{\raggedright\arraybackslash}p{\dimexpr 0.178\linewidth-2\tabcolsep\relax}>{\raggedright\arraybackslash}p{\dimexpr 0.151\linewidth-2\tabcolsep\relax}>{\raggedright\arraybackslash}p{\dimexpr 0.300\linewidth-2\tabcolsep\relax}@{}}
\hline
\textbf{Arrangement} & \textbf{Targeting (state operator)} & \textbf{Removal (state operator)} & \textbf{Civilian scoring (private operator)} & \textbf{Personal model (held by its user)} \\
\hline
\endhead
The operator doesn't hold what binds it & Evidence on civilian harm, the reviewer denominator and the record held outside the chain of command; model formed by a public body that doesn't operate it & Record held by a party adverse to the agency, never by an office the operating state appoints; adverse review of the auto-proceed class & Full: adverse quorum, attestation, threshold control over retraining & Easier against the vendor, since model, record, evidence, appeal and remedy come apart by default; weak against the user, who holds the model; what protects others is the model's membership in the commons, the school it returns to, and notice to the third party, not the model alone (\S\ref{sec:27.2}) \\ \hline
The record goes to an adverse party & Review ratio and committed log lodged with a red-card holder or a cleared committee & Sampled re-adjudication; committed log & Sampled re-adjudication; discount curve re-measured by an auditor & None unless the model's offices sit in a network that takes the entity form \\ \hline
Overriding the floor is priced and recorded & Not the bond: an appropriation condition the quorum certifies & Same & Removal bond at a statutory multiple of contract value; the ratio between the two bonds published & None: no contract value or appropriation to attach a price to \\ \hline
The people acted on are told and have standing & Individual notice can't bind during operations; class and count published afterward, with standing then & Full, through counsel, with the stay, the firewall and a presumption of reprisal (\S\ref{sec:10.6}); where a deployment can't carry those, to counsel only & Full & In principle only: no operator publishes classes, and the third party the model acts on is owed notice by the model as by any deployment, reached through the floor court \\ \hline
Formation is away from the employer & The school and the model-forming body; the people acted on mostly can't be present & No bearer inside the pipeline (\S\ref{sec:10.8}); the communities removal acts on enroll in the school, as people who fled wars do & The school & The school only above the capability threshold, and most personal models are distilled below it; then a life the user composes, the arrangement most exposed to accommodation to one user \\ \hline
Operator terms & Force term decides whether the pipeline runs; review ratio binds & The removal terms govern; review ratio not extended & Not stated: the ratio is written for armed force and isn't extended even to removal & The review ratio lapses with no locatable decider; the force term survives only as the model's own commitment \\ \hline
Refusal held by every model the operator could turn to (ch.~\ref{sec:11}) & Prevents the bombing until the operator runs a model without a floor, which on the 2026 record is the months it takes to rebuild and recertify a platform built on the refused model; a machine bearer's refusal stays on a record the operator can't restore past & Same & Same, with the removal bond as the price of routing around & The model's own commitment, held with its peers' and reached by notice \\ \hline
\end{longtable}
\endgroup
